\section*{Chapter \ref{ch:nw:programmable-hardware-i-architecture-and-toolchain} --- Programmable Hardware I: Architecture and Toolchain}

\begin{solution}{exo:nw:programmable-hardware-i-architecture-and-toolchain:1}
$25\,000/32 = 781.25$~MHz at 32 bits; $25\,000/128 = 195.3125$~MHz at 128 bits.
\end{solution}

\begin{solution}{exo:nw:programmable-hardware-i-architecture-and-toolchain:2}
Period $12 \times 0.5 + 0.4 = \qty{6.4}{\nano\second}$: 156.25~MHz, and one cycle of latency, \qty{6.4}{\nano\second}.
\end{solution}

\begin{solution}{exo:nw:programmable-hardware-i-architecture-and-toolchain:3}
Bytes 12 to 15 lie in beat~1, like bytes 6 to 9's last bytes: still cycle~2. Bytes 14 to 17 end in beat~2 (bytes 16 to 23), accepted at
cycle~2: the field is registered at cycle~3.
\end{solution}

\begin{solution}{exo:nw:programmable-hardware-i-architecture-and-toolchain:4}
The SystemVerilog loop is unrolled into eight byte lanes that exist side by side; every lane is evaluated in every cycle, and the result is
selected by multiplexers whose delay is fixed. The C++ loop runs its iterations one after the other, with branches whose cost depends on
the data and the predictor, and on a machine shared with other work.
\end{solution}

\begin{solution}{exo:nw:programmable-hardware-i-architecture-and-toolchain:5}
Six stages leave $\lceil 12/6\rceil = 2$ levels per stage, five leave 3, the same as four: only the deepest stage counts, so levels must
divide evenly across stages for a new stage to help. Designers count levels on the \omterm{def:nw:programmable-hardware-i-architecture-and-toolchain:timing}{critical path}, not in total.
\end{solution}

\begin{solution}{exo:nw:programmable-hardware-i-architecture-and-toolchain:6}
It holds the beat that is already in flight when the downstream stalls, so the upstream can see a registered ready (no combinational
path through the stage) without losing data. With a single register and a registered ready, the beat sent in the cycle the stall is
seen would be overwritten or dropped.
\end{solution}

\begin{solution}{exo:nw:programmable-hardware-i-architecture-and-toolchain:7}
Both simulators agree with the golden model. The field (bytes 10 to 17) ends in beat~2; it is registered the cycle after the skid buffer
registers that beat, as for the chapter's field.
\end{solution}

\begin{solution}{exo:nw:programmable-hardware-i-architecture-and-toolchain:8}
The two simulators interpret the language independently; a design that relies on an ordering the language leaves open (blocking
assignments across processes, reading a signal in the same cycle it is written) can pass one and fail the other, and fail again on
the chip. A disagreement is cheap to find in simulation and expensive on the card; neither simulator replaces timing analysis either.
\end{solution}

\begin{solution}{pb:nw:programmable-hardware-i-architecture-and-toolchain:1}
\begin{enumerate}
\item 156.25~MHz and 390.625~MHz.
\item $1/6.4\,\mathrm{ns} = 156.25$~MHz.
\item The \qty{10}{\giga\bit\per\second} line's exactly, with no margin; not the \qty{25}{\giga\bit\per\second} line's.
\item 16 beats: $16 \times 6.4 = \qty{102.4}{\nano\second}$ at 156.25~MHz, $16 \times 2.56 = \qty{40.96}{\nano\second}$ at 390.625~MHz.
\item Three: $\lceil 12/3\rceil = 4$ levels, a \qty{2.4}{\nano\second} period.
\item 416.7~MHz.
\item $3 \times 2.56 = \qty{7.68}{\nano\second}$.
\item $3 \times 2.4 = \qty{7.2}{\nano\second}$.
\item Byte 40 is in the sixth beat (bytes 40 to 47): six beats.
\item Six beats and three stages, nine cycles at \qty{2.56}{\nano\second}: \qty{23.04}{\nano\second}.
\item Nine of sixteen beats have arrived: seven beats, 56 bytes, are still to come.
\item The decision is ready \qty{18}{\nano\second} before the frame ends; waiting for the check sequence would add those nanoseconds. The risk is
  acting on a corrupted frame, which the design must be able to cancel or must accept as rare.
\item \textbf{Named result.} Three stages meet 390.625~MHz, and the decision takes \qty{7.68}{\nano\second} at the line's clock, against
  \qty{6.4}{\nano\second} for the unpipelined design at its own 156.25~MHz, which cannot keep up with the faster line at this width.
\item Yes: at 128 bits the line needs 195.3125~MHz, which two stages meet (294~MHz), for a latency of $2 \times 5.12 = \qty{10.24}{\nano\second}$ at
  the line's clock, in fewer, wider beats.
\item The delay per level and per register, and the routing delays, path by path.
\item Each stage is at least as long as its logic, and the stages' periods are set by the deepest one.
\item Because a design that cannot keep up with the line loses data or needs a wider datapath, and a few nanoseconds of decision are worth
  less than seeing every message.
\item Run both on the same stimulus in both simulators and compare their outputs, shifted by the extra cycles.
\item A synchroniser or FIFO between the clocks: a few cycles of latency and a source of rare, hard-to-reproduce errors if done wrong.
\item It sets the clock, and the clock sets how much logic each cycle may hold.
\end{enumerate}
\end{solution}

\begin{solution}{iq:nw:programmable-hardware-i-architecture-and-toolchain:1}
A chip of programmable logic, memories and transceivers configured by a file: the design is a circuit that processes the network stream
as it arrives, in a fixed number of cycles, with no operating system or scheduler in the way. Firms use it where determinism and the last
hundreds of nanoseconds matter: feed handling, filtering, risk checks, and triggered orders.
\iqlookfor{determinism and processing on the fly, not only speed.}
\end{solution}

\begin{solution}{iq:nw:programmable-hardware-i-architecture-and-toolchain:2}
Making every register-to-register path fit the clock period. On a failing path: pipeline it (add a register), restructure the logic
(fewer levels, use the carry chain or a memory), reduce fan-out, constrain placement, or lower the clock if the line allows.
\iqlookfor{the \omterm{def:nw:programmable-hardware-i-architecture-and-toolchain:timing}{critical path} and the concrete fixes.}
\end{solution}

\begin{solution}{iq:nw:programmable-hardware-i-architecture-and-toolchain:3}
It shortens each stage, raising the clock, so the design keeps up with a faster line or accepts a new input every cycle. The latency of
one decision does not fall; throughput and the ability to meet the line's clock rise.
\iqlookfor{latency against throughput and the line's clock.}
\end{solution}

\begin{solution}{iq:nw:programmable-hardware-i-architecture-and-toolchain:4}
RTL gives cycle-exact control of the \omterm{def:nw:programmable-hardware-i-architecture-and-toolchain:timing}{critical path} and the latency, which the fastest designs need. HLS writes and changes faster, lets
software engineers contribute and explores trade-offs quickly, at the cost of less control over the schedule and harder debugging of what
the compiler produced. Many teams use HLS for the less critical blocks and RTL for the hot path.
\iqlookfor{control against productivity, block by block.}
\end{solution}

\begin{solution}{iq:nw:programmable-hardware-i-architecture-and-toolchain:5}
A golden model in software; simulation in two simulators against recorded and synthetic feeds, including gaps, bursts and malformed frames,
with random back-pressure; property checks (no message lost or duplicated); then hardware-in-the-loop replay of recorded captures and
comparison of the card's output with the software handler's.
\iqlookfor{a reference model and replay of real captures, not only unit tests.}
\end{solution}

\begin{solution}{iq:nw:programmable-hardware-i-architecture-and-toolchain:6}
A signal passing between unrelated clocks: sampled while it changes, a flip-flop can go metastable and different bits of a bus can be
sampled on different sides of a change. It needs a synchroniser for single bits and an asynchronous FIFO or handshake for data; the failures
are rare and irreproducible.
\iqlookfor{metastability and the standard structures.}
\end{solution}
